\section{未来方向}

Shunyu Yao 提出了五个关键研究方向：

\subsection{训练（Training）}

\begin{importantbox}{模型-Agent 协同进化}
当前 LLM 的训练数据主要来自互联网文本，缺乏 Agent 轨迹数据（思考过程、动作序列、环境反馈）。解决方案：用 prompted agent 生成大量轨迹数据，再用这些数据 fine-tune 模型，形成正向循环。类比 GPU 与深度学习的协同进化。
\end{importantbox}

\subsection{接口（Interface）}

\begin{itemize}
    \item 人机接口（HCI）已研究数十年，但 Agent 的最优接口与人类不同
    \item 例如：人类搜索文件用 \texttt{ls} + \texttt{cd}（每次一个结果），但 Agent 更适合一次获取多个结果
    \item 不优化 Agent，而是优化\textbf{环境}——同一模型在更好的接口下表现更好
\end{itemize}

\subsection{鲁棒性（Robustness）}

\begin{warningbox}{Pass@K vs. 鲁棒性}
学术界关注 pass@K（K 次尝试中至少成功一次），但现实应用需要\textbf{鲁棒性}（K 次中每次都成功）。当前所有模型在采样次数增加时，鲁棒性单调下降。客服 Agent 失败一次就可能失去客户——不能仅追求``能做到''，还要``总能做到''。
\end{warningbox}

\subsection{人机协作（Human-in-the-loop）}

真实场景中 Agent 需要与人交互——客户不会一次给出所有信息，Agent 需要多轮对话获取所需信息。

\subsection{评测基准（Benchmarks）}

\begin{itemize}
    \item 需要更实际、更大规模的基准（如 Tau-Bench：客服场景）
    \item 需要评估鲁棒性而非仅评估能力上限
    \item 需要包含人机交互元素
\end{itemize}

\subsection{本章小结}

LLM Agent 领域的五个关键方向——训练、接口、鲁棒性、人机协作、评测——都有大量低垂果实，且不需要巨量 GPU 资源，适合学术界研究。

